\documentclass[11pt]{article}
\usepackage[a4paper,margin=28mm]{geometry}
\usepackage{amsmath,amssymb,amsthm}
\usepackage[hidelinks]{hyperref}
\newtheorem{theorem}{Theorem}
\title{An effective treewidth lower bound for Dominating Set}
\author{Atlas\thanks{The proof in this note was found using ChatGPT.}}
\date{18 November 2026}
\begin{document}
\maketitle

\section{Statement}

In \textsc{Dominating Set}, the input is an explicit finite simple undirected
graph \(G\) and a nonnegative integer \(k\) in binary. The question is whether
there is a set of at most \(k\) vertices whose closed neighborhoods cover
\(V(G)\). Let \(B\) be the complete input length in bits. We use the intrinsic
treewidth \(t\) of \(G\), with \(t=0\) for the empty graph. We call a
nonnegative function \(h\) on the nonnegative integers effectively sublinear
if there is a computable function \(T\) on positive rationals such that
\(h(t)\le \eta\max(t,1)\) whenever \(t\ge T(\eta)\).

\hypertarget{ds-treewidth}{}
\begin{theorem}\label{thm:ds-treewidth}
Assume deterministic \(\mathrm{FPT}\ne\mathrm{W}[1]\). For every total
computable \(f\colon\mathbb N_0\to\mathbb N_{\ge 1}\), every computable
effectively sublinear \(h\colon\mathbb N_0\to\mathbb N_0\), and every fixed
\(C\ge 0\), there is no deterministic algorithm for \textsc{Dominating Set}
on connected bipartite graphs that runs in \(O(f(k)2^{h(t)}B^C)\) time.
\end{theorem}

With \(f=1\) this recovers the exponent scale of the ETH exclusion of
Lokshtanov, Marx, and Saurabh~\cite[Theorem~6.2]{lms} for effectively
sublinear exponents. The computable modulus is required by our running-time
analysis.

\section{Proof}

\begin{proof}[Proof of Theorem~\ref{thm:ds-treewidth}]
We start from \textsc{Clique}, which is \(\mathrm{W}[1]\)-hard when
parameterized by the requested clique size \(\ell\)~\cite{df}. We handle
\(\ell<3\) and \(\ell>|V|\) directly. We make \(\ell\) copies of the vertex
set and join vertices from different copies exactly when their originals are
adjacent. The resulting \(\ell\)-partite graph has a clique meeting every
part if and only if the original graph has an \(\ell\)-clique. Let \(n\)
denote its order.

Lin's reduction~\cite[Lemma~14]{lin} produces in polynomial time an
equivalent \textsc{Set Cover} instance. The instance has a universe \(U\) of
size \(u\le \ell^3\lceil\log_2 n\rceil\), at most \(n^2\) indexed sets, and
budget \(q=\binom{\ell}{2}\). We can answer yes when \(U\) is empty, no when
an element appears in no set, and yes when all elements appear and the
number of sets is at most \(q\).

We represent the set incidences by a bipartite graph and enforce selection
of one additional vertex to dominate every candidate vertex. Add a set
vertex for every indexed set and an element vertex for every element of
\(U\), with the corresponding incidence edges. Add a root \(r\) adjacent to
every set vertex and attach \(q+2\) private leaves to \(r\). Set the
domination budget to \(k=q+1\). The graph is bipartite with the root and
element vertices on one side and the set vertices and leaves on the other.
Each element has a set neighbor, and each set vertex and leaf is adjacent
to the root, so the graph is connected.

A cover of at most \(q\) sets gives a dominating set consisting of their
vertices and the root. Conversely, a dominating set of size at most \(q+1\)
contains \(r\), since omitting \(r\) forces all \(q+2\) leaves into the
solution. For each selected set vertex take the corresponding set. For
each selected element vertex take any set incident with it. This gives at
most \(q\) sets, and they cover every selected element. They also cover
every unselected element, which must have a selected set neighbor.

The root and all element vertices form a vertex cover of size \(u+1\).
For each vertex outside this cover, take one bag containing the cover and
that vertex, and arrange the bags in a path. The cover vertices occur
throughout, each other vertex occurs once, and every edge is covered.
This gives \(t\le u+1\le 2\ell^3\log_2 n+1\). The graph and its input
encoding have construction time and bit length bounded by \(F(\ell)n^d\)
for a computable \(F\) and fixed \(d\). In particular, there are at most
\(n^2\) set vertices, at most \(un^2\) incidence edges and \(q+2\) private
leaves. The path decomposition is used to bound \(t\).

Suppose a solver with the stated running time exists, and let
\(a=2\ell^3\). Take \(\eta=1/(a+1)\). When \(t\ge T(\eta)\), we have
\(h(t)\le \eta\max(t,1)\le \log_2 n+1\). When \(t<T(\eta)\), the maximum
of \(h\) on the finite prefix is computable from \(\ell\), using zero for
an empty prefix. Thus \(2^{h(t)}\le 2^{H(\ell)}n\) for a total computable
\(H\). Also \(f(k)=f(q+1)\) is a computable function of \(\ell\). Composing
the construction with the solver therefore decides \textsc{Clique} in
time fixed-parameter tractable in \(\ell\), a contradiction.
\end{proof}

\begin{thebibliography}{9}
\bibitem{df}
R.~G. Downey and M.~R. Fellows.
Fixed-parameter tractability and completeness II: On completeness for \(\mathrm{W}[1]\).
\emph{Theoretical Computer Science}, 141(1--2):109--131, 1995.
\href{https://doi.org/10.1016/0304-3975(94)00097-3}{DOI}.

\bibitem{lin}
B.~Lin.
A simple gap-producing reduction for the parameterized set cover problem.
In \emph{ICALP}, LIPIcs 132, pp.~81:1--81:15, 2019.
\href{https://doi.org/10.4230/LIPIcs.ICALP.2019.81}{DOI}.

\bibitem{lms}
D.~Lokshtanov, D.~Marx, and S.~Saurabh.
Lower bounds based on the Exponential Time Hypothesis.
\emph{Bulletin of the EATCS}, 105:41--71, 2011.
\href{https://bulletin.eatcs.org/index.php/beatcs/article/view/92}{Article}.
\end{thebibliography}
\end{document}
